% ==============================================================
\section{提案手法}\label{sec:method}
% ==============================================================

(1)反密集区間検出，(2)コントラスト密集パターンマイニング，
(3)構造変化の統計的有意性検定のアルゴリズムを提示する．

\subsection{反密集区間検出}\label{sec:anti_dense_algo}

アルゴリズム~\ref{alg:anti_dense}は，密集区間検出の閾値交差論理を
逆転させることで反密集区間を計算する．
核心的な知見は，同じスライディングウィンドウカウント機構を再利用でき，
比較演算子と区間追跡ロジックのみを変更すればよいという点である．

\begin{algorithm}[t]
\caption{反密集区間検出}
\label{alg:anti_dense}
\begin{algorithmic}[1]
\REQUIRE ソート済みタイムスタンプ$\mathbf{t} = [t_1, \ldots, t_n]$，
  ウィンドウサイズ$W$，低閾値$\theta_{\text{low}}$
\ENSURE 反密集区間の集合
  $\overline{\mathcal{D}}(P, W, \theta_{\text{low}})$
\STATE $\mathcal{A} \leftarrow \emptyset$ \COMMENT{結果集合}
\STATE 臨界位置を計算：
  $C = \{t_i\} \cup \{t_i - W\} \cup \{t_i + 1\} \cup \{0, t_n\}$
\STATE $C$をソートし$[0, t_n]$の範囲にフィルタ
\STATE $\texttt{in\_anti} \leftarrow \texttt{false}$;
  $\texttt{start} \leftarrow \texttt{null}$
\FOR{$C$の各要素$c$を順に}
  \STATE $\texttt{count} \leftarrow |\{t_j \in \mathbf{t} \mid
    c \leq t_j \leq c + W - 1\}|$ \COMMENT{二分探索}
  \IF{$\texttt{count} < \theta_{\text{low}}$}
    \IF{$\texttt{in\_anti}$でない}
      \STATE $\texttt{in\_anti} \leftarrow \texttt{true}$;
        $\texttt{start} \leftarrow c$
    \ENDIF
  \ELSE
    \IF{$\texttt{in\_anti}$}
      \STATE $\mathcal{A} \leftarrow \mathcal{A} \cup
        \{[\texttt{start}, c-1]\}$
      \STATE $\texttt{in\_anti} \leftarrow \texttt{false}$
    \ENDIF
  \ENDIF
\ENDFOR
\IF{$\texttt{in\_anti}$}
  \STATE $\mathcal{A} \leftarrow \mathcal{A} \cup
    \{[\texttt{start}, t_n]\}$
\ENDIF
\RETURN $\mathcal{A}$
\end{algorithmic}
\end{algorithm}

\paragraph{計算量}
アルゴリズムは$O(n)$個の臨界位置でサポートを評価し，各評価には
二分探索による$O(\log n)$時間を要する．
総計算量は$O(n \log n)$であり，密集区間検出アルゴリズムと一致する．

\paragraph{枝刈りのための逆単調性}
定理~\ref{thm:anti_monotone}により，上位集合$Q \supset P$が
候補領域に反密集区間を持たなければ，$P$もそこに反密集区間を持たない
（$s_P(t) \geq s_Q(t)$であるため）．
これは\emph{トップダウン}マイニング戦略を可能にする：
大きなアイテムセットから開始し，上位集合が関心領域に反密集区間を
持たない部分集合を枝刈りする．

\subsection{コントラスト密集パターンマイニング}\label{sec:contrast_algo}

トランザクション列を二つのレジーム$R_1 = [0, \tau]$と$R_2 = [\tau+1, N-1]$に
分割するレジーム境界$\tau$が与えられたとき，
アルゴリズム~\ref{alg:contrast}は各パターンのトポロジー変化を分類する．

\begin{algorithm}[t]
\caption{コントラスト密集パターンマイニング}
\label{alg:contrast}
\begin{algorithmic}[1]
\REQUIRE パターン集合$\mathcal{F}$（タイムスタンプ付き），ウィンドウ$W$，
  閾値$\theta$，境界$\tau$，マージン$\delta$
\ENSURE トポロジータイプ付きコントラスト密集パターン$\mathcal{C}$
\STATE $\mathcal{C} \leftarrow \emptyset$
\FOR{各$P \in \mathcal{F}$}
  \STATE $\mathcal{D}_1 \leftarrow$ \textsc{DenseIntervals}$(P, W,
    \theta)$を$R_1$にクリップ
  \STATE $\mathcal{D}_2 \leftarrow$ \textsc{DenseIntervals}$(P, W,
    \theta)$を$R_2$にクリップ
  \STATE $c_1 \leftarrow \text{cov}(\mathcal{D}_1, |R_1|)$;
    $c_2 \leftarrow \text{cov}(\mathcal{D}_2, |R_2|)$
  \IF{$\mathcal{D}_1 = \emptyset \land \mathcal{D}_2 \neq \emptyset$}
    \STATE type $\leftarrow$ \textsc{Emergence}
  \ELSIF{$\mathcal{D}_1 \neq \emptyset \land \mathcal{D}_2 = \emptyset$}
    \STATE type $\leftarrow$ \textsc{Vanishing}
  \ELSIF{$c_2 > c_1 + \delta$}
    \STATE type $\leftarrow$ \textsc{Amplification}
  \ELSIF{$c_1 > c_2 + \delta$}
    \STATE type $\leftarrow$ \textsc{Contraction}
  \ELSE
    \STATE type $\leftarrow$ \textsc{Stable}; \textbf{continue}
  \ENDIF
  \STATE $\mathcal{C} \leftarrow \mathcal{C} \cup \{(P, \text{type})\}$
\ENDFOR
\RETURN $\mathcal{C}$
\end{algorithmic}
\end{algorithm}

\subsection{統計的有意性検定}\label{sec:permtest_algo}

パターンの構造変化が統計的に有意であるかを評価するため，
定義~\ref{def:permtest}で記述した置換検定を用いる．
手順は以下のとおりである：

\begin{enumerate}
\item 観測コントラスト統計量
  $\Delta_{\text{obs}} = \text{cov}_2(P) - \text{cov}_1(P)$を計算する．
\item $b = 1, \ldots, B$に対して：
  \begin{enumerate}
  \item タイムスタンプ位置のランダム置換を生成する
    （$[0, N-1]$からの非復元抽出）．
  \item 密集区間とコントラスト統計量$\Delta^{(b)}$を再計算する．
  \end{enumerate}
\item $p$値を計算する：
  $p = (1 + \sum_{b} \mathbf{1}[|\Delta^{(b)}| \geq |\Delta_{\text{obs}}|]) / (1 + B)$．
\item すべての検定対象パターンに対して
  Benjamini--Hochberg~\cite{benjamini1995fdr}補正を適用する．
\end{enumerate}

\paragraph{計算量}
各置換には密集区間計算のための$O(n \log n)$時間を要する．
$B$回の置換と$|\mathcal{F}|$個のパターンに対し，
総コストは$O(|\mathcal{F}| \cdot B \cdot n \log n)$である．
実用的には$B = 199$または$999$で信頼性のある$p$値推定に十分である．

\subsection{密集区間マイニングとの統合}

本研究の反密集アルゴリズムとコントラストアルゴリズムは，
既存の密集区間マイニングパイプラインにシームレスに統合される：
\begin{enumerate}
\item 標準的なAprioriベースの密集区間マイニングを実行し，
  $\mathcal{F}$とその密集区間を取得する．
\item $\mathcal{F}$の各パターンに対し，
  アルゴリズム~\ref{alg:anti_dense}を用いて反密集区間を計算する．
\item レジーム境界が与えられた場合，アルゴリズム~\ref{alg:contrast}を用いて
  トポロジー変化を分類し，有意性を検定する．
\end{enumerate}
コアマイニングアルゴリズムの変更は不要である．
